% Generated draft: requires mathematical and source review.
\documentclass[11pt]{article}
\usepackage[margin=25mm]{geometry}
\usepackage{fontspec}
\setmainfont{Latin Modern Roman}
\usepackage{amsmath,amssymb,mathtools,unicode-math}
\setmathfont{Latin Modern Math}
\usepackage{enumitem,needspace,xurl,hyperref}
\setlength{\parindent}{0pt}
\setlength{\parskip}{5pt}
\setlength{\emergencystretch}{3em}
\newcommand{\N}{\mathbb{N}}
\newcommand{\Z}{\mathbb{Z}}
\newcommand{\Q}{\mathbb{Q}}
\newcommand{\R}{\mathbb{R}}
\newcommand{\C}{\mathbb{C}}
\newcommand{\F}{\mathbb{F}}
\newcommand{\A}{\mathbb{A}}
\newcommand{\PP}{\mathbb{P}}
\newcommand{\E}{\mathbb{E}}
\newcommand{\eps}{\varepsilon}
\newcommand{\dd}{\,\mathrm{d}}
\begin{document}
\section*{UnsolvedMath ID: 30004473}
\textbf{Source problem code:} catalogue entry
\subsection{Definitions and mathematical statement}
\textbf{Formulation and unresolved data.} The supplied record uses a reductive triple $(G,B,T)$, a Whittaker datum $\psi$, a category $\operatorname{Rep}_{[T,1]}G$, an algebra $Z_G$, a dual group $\check G$, and a space $\mathcal X_{\check G}$. It does not specify the ground field, the representation coefficients, the precise class of groups or representations, the definition of the subscript $[T,1]$, or the construction of $\mathcal X_{\check G}$. Thus the following is a conditional statement of the supplied conjectural target, rather than a completely specified conjecture.

Write
\[
\mathcal C_G=D^+(\operatorname{Rep}_{[T,1]}G),
\qquad
\mathcal D_G=D^+\operatorname{QCoh}([\mathcal X_{\check G}/\check G]).
\]
Here $D^+$ denotes bounded-below derived-category notation, $\operatorname{QCoh}$ denotes quasi-coherent sheaves, and $[\mathcal X_{\check G}/\check G]$ denotes the indicated quotient stack, once the space and group action have been specified. The precise derived-category conventions remain unspecified.

\textbf{Conjectural target as supplied.} For every reductive triple $(G,B,T)$ with Whittaker datum $\psi$ in the intended, unspecified setting, does there exist a fully faithful $Z_G$-linear functor
\[
R_G^\psi:\mathcal C_G\longrightarrow\mathcal D_G
\]
with all of the following properties?
\begin{enumerate}
\item For every parabolic subgroup $P$ with $B\subseteq P\subseteq G$, the functor is compatible with parabolic induction.
\item When $G$ is a torus, the functor agrees with the canonical equivalence referred to in the source record.
\item The compactly induced Whittaker object is sent to the structure sheaf:
\[
R_G^\psi(W_G^\psi)\simeq
\mathcal O_{[\mathcal X_{\check G}/\check G]},
\]
where $W_G^\psi$ is only a label here for the object mentioned in the record; its construction and placement in $\mathcal C_G$ are not supplied.
\end{enumerate}
Full faithfulness means that, for every $A,B\in\mathcal C_G$, the induced map
\[
\operatorname{Hom}_{\mathcal C_G}(A,B)
\longrightarrow
\operatorname{Hom}_{\mathcal D_G}(R_G^\psi A,R_G^\psi B)
\]
is a bijection. The requirement of $Z_G$-linearity presupposes specified compatible actions of $Z_G$ on the two categories, which the record does not provide. Likewise, the parabolic compatibility cannot be made into a precise comparison isomorphism without the relevant functors, their normalizations, and any coherence requirements. The torus equivalence and the compact-induction normalization also require source definitions absent from the supplied record. These omissions must be resolved before this schematic target can be treated as a precise mathematical statement.
\subsection{Short English statement}
The supplied conjecture asks for fully faithful $Z_G$-linear functors from the indicated representation categories to quasi-coherent sheaves on dual-group quotient stacks, compatible with every parabolic induction and normalized by the torus equivalence and the Whittaker object. Essential definitions and comparison functors are missing, so this is a statement-only draft requiring mathematical review.
\subsection{Sources}
\begin{itemize}
\item Local UnsolvedMath registry: supplied complete source record, permanent ID $30004473$, title ``Whittaker-Normalized Functors for Reductive Groups,'' source problem code \texttt{catalogue entry}; collection: Oberwolfach Reports Open Problems. No individual proposer is supplied; the proposed year is recorded as 2020.
\item Supplied report citation: \emph{Arithmetic Geometry (2021)}, Oberwolfach Reports, pp.\ 1023--1082, DOI \texttt{10.4171/owr/2020/20}. Source URL and literature-source link labeled ``original OWR report'': \url{https://doi.org/10.4171/owr/2020/20}. The record describes Oberwolfach Reports as documenting mathematical workshops and their open-problem discussions. This report citation is supplied provenance and has not been independently consulted here.
\item Supplied literature-source link labeled ``Iwahori--Hecke derived category'': \url{https://doi.org/10.1112/S0010437X23007145}. No authors, theorem identifiers, or precise correspondence with the stated conjecture are supplied.
\item Literature assessment dated exactly \texttt{2026-08-22}: the supplied curation labels the problem ``partially solved'' and states that the fully faithful Whittaker-normalized functor and its parabolic compatibilities remain conjectural in general, while principal-block and related cases have substantial proved results. This is supplied dated curation, not an independently verified assessment. The background explicitly requests independent verification of time-sensitive status and cited claims.
\item The registry labels the statement \texttt{corrected\_verified} and reports that it was checked against the official Oberwolfach report. These are supplied verification claims, not verification performed for this draft.
\end{itemize}
\end{document}
